\documentclass[11pt]{article}
\usepackage[utf8]{inputenc}
\usepackage[T1]{fontenc}
\usepackage{geometry}
\geometry{a4paper, margin=2.5cm}

\title{Credit Platform - Rapport Technique}
\author{Equipe Dev}
\date{\today}

\begin{document}
\maketitle
\tableofcontents
\newpage

\section{Introduction}
\subsection{Problematique}
Gestion complexe des demandes de credit multi-canaux (Email, WhatsApp).
\subsection{Solution}
Credit Platform centralise et analyse automatiquement les demandes via IA et Qdrant.

\section{Architecture}
5 Phases :
\begin{enumerate}
    \item Ingestion (Gmail, WhatsApp)
    \item NLP (Nettoyage, Extraction)
    \item Vectorisation (Qdrant Cloud)
    \item Decision Metier
    \item Dashboard
\end{enumerate}

\section{Integration Qdrant}
Recherche par similarite (pas de classification).
Collection: email\_memory, Vecteur: 384 dimensions, Distance: Cosine.

\section{Status du Projet}
\subsection{Fait (Work Done)}
\begin{itemize}
    \item Scripts d'ingestion (Gmail/WhatsApp/Banque) operationnels
    \item Traitement NLP et nettoyage en place
\end{itemize}

\subsection{A Faire (To Do)}
\begin{itemize}
    \item \textbf{Setup Qdrant} : Configuration et connexion Infrastructure
    \item \textbf{Dataset Synthétique} : Génération des emails de référence
    \item Vectorisation effective des messages
    \item Moteur de recherche de similarite
    \item Workflows de decision
    \item Interface utilisateur
\end{itemize}

\end{document}
